% Part: computability
% Chapter: computability-theory
% Section: prop-reduce

\documentclass[../../../include/open-logic-section]{subfiles}

\begin{document}

\olfileid{cmp}{thy}{ppr}
\olsection{न्यूनीकरणक्षमतेचे गुणधर्म}

$A$ चे~$B$ कडे न्यूनीकरण होत असेल, तर आपण $A \leq_m B$
असे लिहितो. या चिन्हांकनावरून असे सुचते की $A \leq_m B$
असल्यास $A$ हे~$B$ पेक्षा ``अधिक कठीण नाही'' आणि $B$
हे~$A$ इतके किंवा त्याहून ``अधिक कठीण'' आहे. तसेच नेहमीचे
संख्यांवरील $\le$ जसे क्रम लावते, तसे $\le_m$ संचांना
(किंवा निर्णय समस्यांना) त्यांच्या गुंतागुंतीनुसार क्रम लावते.
पुढील दोन विधाने ही अंतःप्रेरणा समर्थित करतात. पहिले विधान
$\le_m$ संक्रमक आहे असे सांगते.

\begin{prop}
\ollabel{prop:trans-red}
$A \leq_m B$ आणि $B \leq_m C$ असल्यास $A \leq_m C$.
\end{prop}

\begin{proof}
$A$ चे~$B$ कडे न्यूनीकरण आणि $B$ चे~$C$ कडे न्यूनीकरण
यांचे संयोजन केल्यास $A$ चे~$C$ कडे न्यूनीकरण मिळते.
\end{proof}

\begin{prob}
$f$ आणि~$g$ ही अनुक्रमे $A$ चे~$B$ कडे आणि $B$ चे~$C$
कडे अनेक-एक न्यूनीकरणे असतील, तर $\comp{f}{g}$ हे $A$
चे~$C$ कडे अनेक-एक न्यूनीकरण आहे, हे दाखवून
\olref{prop:trans-red} सिद्ध करा.
\end{prob}

\begin{prop}
\ollabel{prop:reduce}
$A$ आणि $B$ हे कोणतेही संच असू द्या आणि $A \leq_m B$
असे समजा.
\begin{enumerate}
\item $B$ संगणनक्षमपणे प्रगणनीय असेल, तर $A$ देखील तसाच आहे.
\item $B$ संगणनक्षम असेल, तर $A$ देखील तसाच आहे.
\end{enumerate}
\end{prop}

\begin{proof}
$f$ हे $A$ चे~$B$ कडे अनेक-एक न्यूनीकरण असू द्या.
पहिल्या विधानासाठी, $B$ हे आंशिक फलन~$g$ चे परिभाषाक्षेत्र
असेल, तर $A$ हे~$\comp{f}{g}$ चे परिभाषाक्षेत्र आहे हे
तपासा:
\begin{align*}
x \in A & \text{ तेव्हाच } f(x) \in B \\
& \text{ तेव्हाच }  g(f(x)) \fdefined.
\end{align*}

दुसऱ्या विधानासाठी, $B$ संगणनक्षम असेल, तर $B$
आणि~$\Complement{B}$ हे दोन्ही संगणनक्षमपणे प्रगणनीय
असतात (\olref[cmp]{thm:ce-comp}), हे आठवा. $f$ हे
$\Complement{A}$ चे $\Complement{B}$ कडेही अनेक-एक
न्यूनीकरण आहे, हे तपासणे कठीण नाही. त्यामुळे या सिद्धतेच्या
पहिल्या भागानुसार $A$ आणि~$\Complement{A}$ हे
!!{computably
enumerable} आहेत. म्हणून \olref[cmp]{thm:ce-comp} मुळे
$A$ देखील संगणनक्षम आहे. (पर्यायाने,
$\Char{A} = \comp{f}{\Char{B}}$ हे तपासा; मग
$\Char{B}$ संगणनक्षम असेल, तर~$\Char{A}$ देखील आहे.)
\end{proof}

\begin{prob}
$f$ हे $A$ चे~$B$ कडे अनेक-एक न्यूनीकरण आहे असे
समजा. मग $f$ हे $\Complement{A}$ चे $\Complement{B}$
कडेही अनेक-एक न्यूनीकरण आहे हे दाखवा.
\end{prob}

\begin{prob}
$f\colon \Nat \to \Nat$ हे $A$ चे~$B$ कडे अनेक-एक न्यूनीकरण असेल, तर
$\Char{A}
= \comp{f}{\Char{B}}$ हे दाखवा.
%READERNOTE{OLCMP-022 — गोठवलेल्या इंग्रजी प्रश्नात अनेक-एक न्यूनीकरण f चे प्रकारलेखन f:A→B असे आहे; मात्र आधीची व्याख्या f:N→N अशी आहे आणि जातिबोधक फलांचे समीकरण संपूर्ण N वर तपासायचे आहे. म्हणून येथे आधीच्या व्याख्येशी सुसंगत प्रकारलेखन वापरले आहे.}
\end{prob}

\begin{digress}
\emph{ट्यूरिंग न्यूनीकरणक्षमता} ही न्यूनीकरणक्षमतेची अधिक
सर्वसाधारण संकल्पना इतर संदर्भांत, विशेषतः अनिर्णेयतेचे निष्कर्ष
सिद्ध करण्यासाठी, उपयोगी आहे. \olref[cmp]{cor:comp-k} नुसार
$K_0$ चा पूरक संच~$K_0$ कडे न्यूनीत करता येत नाही, कारण
तो संगणनक्षमपणे प्रगणनीय नाही. पण अंतःप्रेरणेने, $K_0$
विषयीच्या प्रश्नांची उत्तरे माहीत असतील, तर त्याच्या पूरकाविषयीची
उत्तरेही माहीत असतील. संगणनक्षम प्रक्रिया~$B$ विषयी प्रश्न
विचारून~$A$ विषयीच्या प्रश्नांची उत्तरे ठरवू शकत असेल, तर
$A$ हे~$B$ कडे ट्यूरिंग न्यूनीत करता येते असे म्हणतात.
अनेक-एक न्यूनीकरणक्षमतेपेक्षा ही अधिक व्यापक संकल्पना आहे:
अनेक-एक न्यूनीकरणात (1)~$B$ विषयी एकच प्रश्न विचारता येतो,
आणि (2)~``होय'' या उत्तराचे $A$ विषयीच्या प्रश्नासाठी
``होय'' असेच, तसेच ``नाही'' चे ``नाही'' असेच रूपांतर व्हावे
लागते. $A$ हे~$B$ कडे ट्यूरिंग न्यूनीत करता येत असेल
आणि $B$ संगणनक्षम असेल, तर $A$ देखील संगणनक्षम असतो;
परंतु आपण पाहिल्याप्रमाणे, संगणनक्षमपणे प्रगणनीय असण्याबाबत
असेच विधान खरे ठरत नाही.

आपण चर्चिलेल्या न्यूनीकरणक्षमतेच्या वेगवेगळ्या संकल्पनांचा
आणि त्यांतील फरकांचा विचार करा. मात्र या प्रकरणात आपण फक्त
अनेक-एक न्यूनीकरणक्षमता वापरू. आधीच्या परिच्छेदातील दोन्ही
प्रकारांना संगणकीय गुंतागुंतीत समांतर संकल्पना आहेत; त्यांत
ट्यूरिंग यंत्रे बहुपदी वेळेत चालण्याची अतिरिक्त अट असते.
अनेक-एक न्यूनीकरणक्षमतेचे गुंतागुंत-संबंधी रूप
\emph{कार्प न्यूनीकरणक्षमता}, तर ट्यूरिंग न्यूनीकरणक्षमतेचे
तसे रूप \emph{कुक न्यूनीकरणक्षमता} म्हणून ओळखले जाते.
\end{digress}

\end{document}
